% ═════════════════════════════════════════════════════════════════════════════
\section{Cutoff Methods and Dynamic Cutoffs}
% ═════════════════════════════════════════════════════════════════════════════

One of the most elegant ideas in parameterized verification is the
\emph{cutoff theorem}. The key insight is that, for certain classes of
systems, one does not need to verify infinitely many instances. Instead, it
suffices to verify systems up to some finite bound, the cutoff, and
correctness for all larger instances follows automatically~\cite{EN95,EK03}.

Intuitively, this works because processes in many symmetric systems are
interchangeable. If the behavior of a system with~$n$ processes can be
faithfully represented by a system with at most~$k$ processes (the cutoff),
then verifying the $k$-process system is enough. Cutoff results have been
established for important classes such as token-passing rings, broadcast
protocols, and systems where processes communicate in restricted ways. A
significant body of work has developed cutoffs for mutual exclusion and cache
coherence protocols~\cite{EK03}.

However, classical cutoff theorems are \emph{static}: the bound~$k$ is fixed
ahead of time and must be established before verification begins. This can be
limiting: the cutoff may be large, difficult to compute, or may not exist
at all for systems with richer interaction patterns or more complex topologies.

\emph{Dynamic cutoff methods} address this limitation by allowing the
abstraction boundary to be determined adaptively, during the verification
process itself, rather than fixed in advance. A prominent and well-developed
technique in this family is \textbf{view abstraction}~\cite{AHH13,AHH16}.
Rather than committing to a single small representative system upfront, the
method reasons about what a process can \emph{observe} about the rest of the
system, and uses those local observations to construct an on-the-fly
abstraction.

The starting point is to model the global state of an $n$-process system as a
\emph{configuration}: a sequence of local process states. A \emph{view} of
size~$k$ is a subsequence of length~$k$, a projection of the global
configuration onto a subset of~$k$ processes. Instead of tracking full
configurations (whose length varies with~$n$), the method tracks the set of
all views of size up to~$k$ that are reachable across \emph{any} instance of
the system.

For a given precision parameter $k \in \mathbb{N}$, the framework runs two
procedures in parallel. The first is an \emph{under-approximation}: it
performs concrete reachability analysis on instances of fixed size up to~$k$,
and reports a bug if a bad configuration is found. The second is an
\emph{over-approximation} based on abstract interpretation. It defines an
abstraction function~$\alpha_k$, which maps a configuration to its set of
subwords of size up to~$k$, and a corresponding concretization~$\gamma_k$,
which maps a set of views back to all configurations whose subwords are all
contained in that set. A fixpoint $V_k$ is then computed over views using an
abstract post operator. The key soundness property is that the set of
reachable configurations~$R$ satisfies $R \subseteq \gamma_k(V_k)$ for
all~$k$, and the approximation becomes tighter as $k$ grows:
$\gamma_{k+1}(V_{k+1}) \subseteq \gamma_k(V_k)$.

A crucial implementation insight is that the abstract post operator can be
computed using only concrete configurations of size $k{+}1$, rather than the
full (infinite) concretization. This is because the witnesses required to
enable a global transition consist of a single extra process; accommodating
one additional process beyond the view size is always sufficient. This makes
the fixpoint computation entirely finite and automatic.

If the over-approximation establishes $\gamma_k(V_k) \cap \mathit{Bad} =
\emptyset$, the system is declared safe. If the under-approximation finds a
reachable bad configuration, the system is declared unsafe. In either case,
the current~$k$ is a \emph{dynamic cutoff}, a point discovered during
verification beyond which the search need not continue. If neither procedure
is conclusive, the precision is increased by incrementing~$k$ and the
analysis is repeated.

The procedure is guaranteed to terminate whenever the set of reachable
configurations is \emph{downward closed} with respect to the subword ordering
on configurations. This condition holds in particular for \emph{monotonic
systems}: systems whose transition relation is monotone with respect to the
subword ordering, in the sense that adding processes to a configuration can
only enable more transitions, not fewer. This covers a broad and well-studied
class of systems, including lossy channel systems, Petri nets (where the
method solves the coverability problem), and timed Petri nets. For Burns'
mutual exclusion protocol, for instance, the method terminates at~$k = 2$:
already with two-process views it can be shown that no bad configuration is
reachable for any system size.

View abstraction has been successfully applied to a range of protocols,
including mutual exclusion algorithms such as the fine-grained version of
Szymanski's protocol, a benchmark that had not been verified automatically
by other methods, as well as cache coherence protocols and distributed
systems with complex process interactions. The framework extends naturally to
other process topologies such as multisets, rings, and trees, and to settings
where global transitions are not assumed to be atomic~\cite{AHH13}.

Together, static cutoff theorems and dynamic approaches like view abstraction
form a complementary toolkit: static cutoffs offer clean theoretical
guarantees for well-structured protocol families, while dynamic methods
provide greater flexibility and automation for protocols that fall outside
those families.
